\documentclass[11pt,a4paper]{article}

\usepackage[margin=2.5cm]{geometry}
\usepackage[T1]{fontenc}
\usepackage{lmodern}
\usepackage{microtype}
\usepackage{parskip}
\usepackage{hyperref}

\newcommand{\code}[1]{\texttt{#1}}

\title{IS4302 Writing Assignment 1\\[0.3em]\large Group Project: AgentBL --- AI-Priced eBL-Backed RWA Trade Finance}
\author{Your Name (A0000000X)}
\date{\today}

\begin{document}
\maketitle

\section*{Question 1: Centralized, decentralized, or partially decentralized?}

AgentBL is \emph{partially decentralized}. It lets an exporter pledge an
electronic bill of lading (eBL) and raise cash by selling discounted RWA tokens
backed by the cargo. An AI agent decides the issue price. The parts that must be
trusted by parties who do not trust each other live in smart contracts on the
Injective testnet. \code{EBLRegistry} records who holds title to the cargo, its
endorsement history and whether it is already pledged. It also refuses to
register the same cargo hash twice, which targets double financing (as in the
2014 Qingdao metals fraud, where the same metal was pledged to several banks).
\code{RWAOfferingPool} enforces the offering life cycle (\code{Open} $\to$
\code{Funded} $\to$ \code{InTransit} $\to$ \code{Repaid}, or \code{Paused},
\code{Frozen} or \code{Liquidation}). Neither the exporter nor the platform can
skip a step, and settlement requires both a payment proof and an arrival proof.
Every AI pricing decision is committed to \code{RiskPricingOracle} with an
evidence hash, a quote hash and a replay-protected decision ID, so the price an
investor paid can be audited later. Everything else runs off-chain on our
server: parsing and cross-checking invoices, bills of lading and insurance
policies, cargo valuation, LLM-based risk scoring and investor KYC.

We split the system this way for three reasons. First, LLM inference is costly
and non-deterministic, so validators cannot re-execute it to reach consensus;
only its output and a hash of its inputs belong on-chain. Second, trade
documents reveal counterparties, prices and margins, so publishing them on a
public ledger would leak commercial secrets. Storing only hashes lets a party
later prove what was used without disclosing it. Third, tokenized trade
receivables are an investment product: someone must perform KYC and be legally
accountable. An eBL also has legal effect only when a reliable system issues it
(e.g.\ Singapore's MLETR-based Electronic Transactions Act). Hence an admin key
mints eBLs and manages \code{InvestorComplianceGate}. The price of this design
is that the owner, executor and oracle-updater keys remain points of trust. The
contracts do not remove this trust, but they make every use of it visible and
auditable. A production version would move these keys to a multisig shared by
the carrier, the bank and the platform, and replace the single executor with an
oracle (Question~2).

\section*{Question 2: Does the application need an oracle, and does ASTREA apply?}

Yes. A smart contract cannot observe ships or banks, so AgentBL needs off-chain
information of three kinds:
\begin{enumerate}
  \item \textbf{Physical events}: vessel departure (\code{markInTransit}),
        cargo arrival (\code{recordCargoArrival}), and port closures, typhoons
        and war-risk zones, which can trigger \code{PAUSE} or \code{FREEZE}.
  \item \textbf{Financial facts}: whether the importer has paid
        (\code{recordImporterPayment}, which unlocks settlement), and commodity
        prices such as LME copper.
  \item \textbf{The AI risk price itself}, pushed through
        \code{RiskPricingOracle}.
\end{enumerate}
Today all three are centralized oracles, i.e.\ allow-listed updater and
executor keys. A compromised key, or a platform colluding with an exporter,
could mark a pool as repaid or keep an old high price after a war breaks out.

ASTREA applies, but only to part of the problem. It works for boolean
propositions whose truth is publicly observable, so that randomly assigned
voters and staked certifiers can check them independently and the honest answer
becomes the focal point. Some of our triggers fit this well, for example
``vessel IMO~9xxxxxx berthed at Yangshan before 30~June'' or ``Shanghai port
suspended operations on date~$D$''. Anyone can check these against AIS or
port-authority data, and settlement can tolerate ASTREA's voting latency. Using
ASTREA here would remove our single executor key. It does not fit the rest.
(i)~The AI price is a forecast, not a fact. Whether a price of 0.83 was
``correct'' is known only after maturity, so voters cannot verify it. Reducing
it to ``is risk high?'' discards the information we need, and voters would
simply copy the agent's answer. What we need here is accountability (committed
inputs and evidence hashes compared ex post), not truth discovery.
(ii)~Bank payments and document authenticity are private. Voters could only
check them by seeing confidential documents, which defeats our privacy design.
Signed attestations from liable parties (the issuing bank or the carrier),
possibly with zero-knowledge proofs, fit better. (iii)~During a crisis,
repricing must take minutes, before the next subscription, and continuous
prices need aggregated data feeds rather than yes/no votes. Moreover, ASTREA is
secure only if certifier stakes exceed what an attacker gains, and one copper
cargo is worth about USD~4.5\,million. We would therefore use a hybrid: ASTREA
for public binary events, bank and carrier attestations for private facts, and
an accountable, evidence-committing oracle for AI pricing.

\end{document}
